% Part: proof-theory
% Chapter: proof-search
% Section: search-algorithm

\documentclass[../../../include/open-logic-section]{subfiles}

\begin{document}

\olfileid{pt}{ps}{sal}

\olsection{Algoritma Panelusuran kanggo \Log{G3c}}

Kita njlentrehaké algoritma panelusuran kanggo
\Log{G3c}. Algoritma iki dudu siji-sijiné kang mungkin,
lan uga dudu kang paling éfisièn. Nanging algoritma iki
bener: yèn sekuèn $\Gamma \Sequent \Delta$ nduwèni
!!a{proof}, pungkasane algoritma mandheg lan nemokaké
!!a{proof}. Algoritma iki uga bisa dipercaya: yèn
mandheg tanpa nemokaké !!a{proof}, utawa ora tau
mandheg, $\Gamma \Sequent \Delta$ pancèn ora valid
saka wiwitan.

Sipat wigati kang kudu diduwèni algoritma iki yaiku
jaminan yèn saben !!{formula} ing
$\Gamma \Sequent \Delta$, utawa ing sekuèn kang
diasilaké sajroning panelusuran, bakal ``direduksi'':
dianggep rumus utama ing aturan inferènsi kang
ditrapaké mundur, lan iki kelakon sakehing kaping kang
dibutuhaké. Sipat iki diarani \emph{kaadilan}.

Algoritma lumaku tataran demi tataran. Asil saben
tataran iku wit sekuèn; asil tataran sabanjuré
diasilaké kanthi ngreduksi !!a{formula} ing saben
sekuèn paling dhuwur kang durung dadi aksioma.

Ing tataran~$0$, kita miwiti saka
$\Gamma \Sequent \Delta$. Kanggo njamin kaadilan,
kita mènèhi indeks angka marang !!{formula}s ing
$\Gamma$ lan $\Delta$, kanthi !!{formula}s kang béda
nampa angka kang béda. Tuladhané, kita bisa
ngétung !!{formula}s saka kiwa menyang tengen.
Angka mau ditulis minangka superskrip. Yèn kita
nggoleki !!a{proof} tumrap $!A, !B \Sequent !C$,
asil tataran~$0$ yaiku
$!A^1, !B^2 \Sequent !C^3$. Angka~$n$ paling gedhé
kang katon minangka superskrip ing sekuèn mawa indeks
iki diarani \emph{indeks maksimal} ($3$ ing tuladha iki).

Yèn sekuèn ngemot kuantor, kita uga kudu ngerti tèrma
kang digunakaké nalika ngreduksi
$\lexists[x][!A(x)]$ ing sisih tengen utawa
$\lforall[x][!A(x)]$ ing sisih kiwa. Kita mung perlu
tèrma kang dibangun saka !!{constant}s $\Obj c_0$,
$\Obj c_1$, $\Obj c_2$, \dots, nganggo lambang fungsi
kang katon ing~$\Gamma$ lan $\Delta$. Kita nganggep
himpunan kabèh tèrma iki wis dienumerasi kanthi urutan
tetep, mula bisa nyebut, tuladhané, ``!!{constant}
kapisan kang ora katon ing $\Gamma \Sequent \Delta$.''
Saben sekuèn mawa indeks ing wit asil algoritma
nduwèni himpunan !!{constant}s kang gegandhèngan.
Himpunan !!{constant}s kanggo sekuèn ing tataran~$0$
iku kabèh !!{constant}s kang katon ing sekuèn iku
(yèn ana), utawa $\{\Obj c_0\}$ yèn ora ana
!!{constant}s.

Anggep ing tataran~$n$ algoritma wis ngasilaké wit
sekuèn mawa indeks lan himpunan !!{constant}s kang
gegandhèngan. Ing tataran~$n+1$, tumrap saben
sekuèn paling dhuwur kita tindakaké langkah iki.

Yèn ana !!a{formula}~$!A$ kang katon ing sisih kiwa
lan tengen sekuèn, utawa sisih kiwa ngemot~$\lfalse$,
kita ora nindakaké apa-apa: sekuèn kaya mangkono
nduwèni !!{proof}s ing~\Log{G3c}, mula sekuèn paling
dhuwur iki ora perlu diolah manèh.

Yèn ora ana !!{formula} nonatomik ing sekuèn,
algoritma mandheg tanpa asil. Ora ana bukti tumrap
$\Gamma \Sequent \Delta$.

Yèn ora mangkono, pilih !!{formula} nonatomik~$!A$
kang indeksé paling cilik, yaiku~$i$. Dadi, sekuèn
paling dhuwur iku awujud $!A^i, \Pi \Sequent \Lambda$
utawa $\Pi \Sequent \Lambda, !A^i$. Anggep $k$
indeks maksimal sekuèn iku, lan $C$ himpunan
!!{constant}s kang gegandhèngan karo sekuèn mau.

Kita ``ngreduksi'' $!A^i$, yaiku ngasilaké sekuèn
paling dhuwur anyar miturut aturan inferènsi kang
ditetepaké déning !!{main operator} saka~$!A$ lan
miturut panggonané $!A^i$ ing sisih kiwa utawa tengen.

\begin{enumerate}
  \item $!A \ident !B \land !C$ lan katon ing sisih
  kiwa: tambahana siji sekuèn paling dhuwur anyar ing
  sadhuwuré $(!B \land !C)^i, \Pi \Sequent \Lambda$:
  \[
  \Axiom$!B^{k+1}, !C^{k+2}, \Pi \fCenter \Lambda$
  \RightLabel{\LeftR{\land}}
  \UnaryInf$(!B \land !C)^i, \Pi \fCenter \Lambda$
  \DisplayProof
  \]
  Himpunan !!{constant}s kang gegandhèngan karo
  sekuèn anyar iku~$C$.
  \item $!A \ident !B \land !C$ lan katon ing sisih
  tengen: tambahana loro sekuèn paling dhuwur anyar ing
  sadhuwuré $\Pi \Sequent \Lambda, (!B \land !C)^i$:
  \[
  \Axiom$\Pi \fCenter \Lambda, !B^{k+1}$
  \Axiom$\Pi \fCenter \Lambda, !C^{k+1}$
  \RightLabel{\RightR{\land}}
  \BinaryInf$\Pi \fCenter \Lambda, (!B \land !C)^i$
  \DisplayProof
  \]
  Himpunan !!{constant}s kang gegandhèngan karo
  sekuèn anyar iku~$C$.

  \item Kasus $!A \ident \lnot !B$,
  $!A \ident !B \lor !C$, lan
  $!A \ident !B \lif !C$ padha carané lan
  dipasrahaké dadi latihan.

  \item $!A \ident \lexists[x][!B(x)]$ lan katon ing
  sisih kiwa: tambahana siji sekuèn paling dhuwur anyar
  ing sadhuwuré
  $\lexists[x][!B(x)]^i, \Pi \Sequent \Lambda$:
  \[
  \Axiom$!B(c)^{k+1}, \Pi \fCenter \Lambda$
  \RightLabel{\LeftR{\lexists}}
  \UnaryInf$\lexists[x][!B(x)]^i, \Pi \fCenter \Lambda$
  \DisplayProof
  \]
  Ing kéné $c$ iku !!{constant} kapisan kang ora ana
  ing~$C$. Himpunan !!{constant}s kanggo sekuèn anyar
  iku $C \cup \{c\}$.
  \item $!A \ident \lexists[x][!B(x)]$ lan katon ing
  sisih tengen: tambahana siji sekuèn paling dhuwur
  anyar ing sadhuwuré
  $\Pi \Sequent \Lambda, \lexists[x][!B(x)]^i$:
  \[
  \Axiom$\Pi \fCenter \Lambda, \lexists[x][!B(x)]^{k+1}, !B(t)^{k+2}$
  \RightLabel{\RightR{\lexists}}
  \UnaryInf$\Pi \fCenter \Lambda, \lexists[x][!B(x)]^i$
  \DisplayProof
  \]
  Ing kéné $t$ tèrma kapisan kang mung ngemot
  !!{constant}s ing~$C$ lan durung digunakaké kanggo
  reduksi tengen $\lexists[x][!B(x)]$ ing sangisoré
  sekuèn iki (utawa $\Obj c_0$ yèn kabèh tèrma kaya
  mangkono wis digunakaké). Himpunan !!{constant}s
  kanggo sekuèn anyar iku~$C$.
  \item Kasus $!A \ident \lforall[x][!B(x)]$
  padha carané lan dipasrahaké dadi latihan.
\end{enumerate}

Yèn ing tataran~$n+1$ ora ana sekuèn anyar kang
ditambahaké ing sadhuwuré sekuèn paling dhuwur
endi waé, algoritma mandheg kanthi kasil. Ing kéné
kita wis nemokaké !!a{proof} tumrap
$\Gamma \Sequent \Delta$ saka wit ing tataran~$n+1$
kanthi mbusak kabèh indeks. Kabèh sekuèn paling
dhuwur awujud $!A, \Pi \Sequent \Lambda, !A$ utawa
$\lfalse, \Pi \Sequent \Lambda$. Kabèh inferènsi
iku inferènsi \Log{G3c} kang bener. Tumrap
inferènsi proposisional iki cetha. Gatèkna yèn ing
saben langkah, himpunan~$C$ saka !!{constant}s kang
gegandhèngan karo sekuèn ngemot kabèh !!{constant}s
kang katon ing sekuèn. Mula, ing reduksi kang nganggo
\LeftR{\lexists} lan \RightR{\lforall}, syarat
eigenvariabel kapenuhi: eigenvariabel~$c$ minangka
!!a{constant} kang ora ana ing himpunan~$C$ kanggo
konklusi, mula $c$ uga ora katon ing konklusi.
Kita olèh:

\begin{prop}
Algoritma panelusuran bukti kanggo \Log{G3c} bener:
yèn mandheg kanthi kasil, sekuèn wiwitan
$\Gamma \Sequent \Delta$ nduwèni !!a{proof}.
\end{prop}





\end{document}
